% --- Archon physics formalization source begin ---
% archon:physics
% archon:covers PhyXMiniProblems/problem_phyx_mini_0112.lean
% archon:source-report reports/phyx_mini/problem_phyx_mini_0112.source.json

\chapter{Physics problem phyx\_mini\_0112}
\label{ch:PhyXMiniProblems_problem_phyx_mini_0112}

\paragraph{Problem source.}
\#\# Physical scenario

Focal Length of a Zoom Lens. Figure shows a simple version of a zoom lens. The converging lens has focal length \$f\_1\$ and the diverging lens has focal length \$f\_2 = -|f\_2|\$. The two lenses are separated by a variable distance \$d\$ that is always less than \$f\_1\$. Also, the magnitude of the focal length of the diverging lens satisfies the inequality \$|f\_2| > (f\_1 - d)\$. To determine the effective focal length of the combination lens, consider a bundle of parallel rays of radius \$r\_0\$ entering the converging lens. \$f\_1 = 12.0 \textbackslash{}, 	ext\{cm\}\$, \$f\_2 = -18.0 \textbackslash{}, 	ext\{cm\}\$, and the separation \$d\$ is adjustable between 0 and 4.0 cm.

\#\# Question

What value of \$d\$ gives \$f = 30.0 \textbackslash{}, 	ext\{cm\}\$?

\#\# Answer choices

- A: A: 1.2cm

- B: B: 2.4cm

- C: C: 3.6cm

- D: D: 0.9cm

\#\# Figure caption (auxiliary; use the image as primary evidence)

The image is a diagram of an optical system involving two lenses. Here’s a detailed description:

- **Lenses**: 
  - There are two lenses depicted:
    - The first lens on the left is a convex lens labeled with focal length \textbackslash{}( f\_1 \textbackslash{}).
    - The second lens on the right is a concave lens labeled with focal length \textbackslash{}( f\_2 = -|f\_2| \textbackslash{}).

- **Rays and Object**:
  - A set of parallel rays (purple) are incident from the left towards the convex lens.
  - These rays converge towards a point labeled \textbackslash{}( Q \textbackslash{}) near the first lens.
  - \textbackslash{}( r\_0 \textbackslash{}) represents the initial height of the rays.
  - After passing through the convex lens, the rays are shown converging and then diverging as they pass through the concave lens to form an image \textbackslash{}( I' \textbackslash{}).
  - Dashed lines depict the path of the rays through both lenses.

- **Distances and Notations**:
  - The distance between the two lenses is labeled \textbackslash{}( d \textbackslash{}).
  - The distance from the second lens to the image \textbackslash{}( I' \textbackslash{}) is \textbackslash{}( s'\_2 \textbackslash{}).
  - The total distance from the object to the final image is labeled \textbackslash{}( f \textbackslash{}).
  - \textbackslash{}( r'\_0 \textbackslash{}) represents the height of the rays after passing through the first lens.

This diagram visually represents the interaction of light rays with the lenses, illustrating how the rays converge and diverge to form an image after passing through both lenses.

\#\# Dataset metadata

Geometrical Optics; Multi-Formula Reasoning, Spatial Relation Reasoning

\paragraph{Recorded answer/context.}
A: A: 1.2cm

\paragraph{Figure/image path.}
/root/proposal\_for\_physic/science-mango/phyx\_mini\_run/phyx\_data/test\_image/112.png

\paragraph{Formalization target.}
create a compiling Lean file with sorry bodies at `PhyXMiniProblems/problem\_phyx\_mini\_0112.lean`. The Lean declarations must preserve the physical quantities, dimensions or dimensional roles, figure labels, governing-law hypotheses, and final relation expressed by this problem.
Use Mathlib/Physlib names found through LeanExplore where available. If a physics API is missing, introduce faithful local abstractions rather than scalar placeholder aliases.

\begin{theorem}[Physics formalization target]
\label{thm:physics:phyx_mini_0112:target}
\lean{PhyXMiniProblems.ProblemPhyXMini0112.problem_phyx_mini_0112}
\uses{def:physics:phyx-mini-0112:phyxminiproblems-problemphyxmini0112-zoomlensdiagram, def:physics:phyx-mini-0112:phyxminiproblems-problemphyxmini0112-lensseparationreadout, def:physics:phyx-mini-0112:phyxminiproblems-problemphyxmini0112-matcheszoomlensfigure, def:physics:phyx-mini-0112:phyxminiproblems-problemphyxmini0112-matcheszoomlensproblemdata, def:physics:phyx-mini-0112:phyxminiproblems-problemphyxmini0112-hasphysicalzoomlensconfiguration, def:physics:phyx-mini-0112:phyxminiproblems-problemphyxmini0112-satisfiesparaxialzoomlenslaws, def:physics:phyx-mini-0112:phyxminiproblems-problemphyxmini0112-matchesanswer}
The assigned autoformalize agent should translate this physics problem into Lean declarations in the covered file, with theorem and lemma proof bodies written as `by sorry`.
\end{theorem}
\begin{proof}
This is an autoformalization task, not a proof task. Produce statements that can later be proved by the physics prover without weakening the physical model.
\end{proof}

% --- Archon declaration topology begin ---
\section{Declaration topology}
\begin{definition}[Optical Length]
  \label{def:physics:phyx-mini-0112:phyxminiproblems-problemphyxmini0112-opticallength}
  \lean{PhyXMiniProblems.ProblemPhyXMini0112.OpticalLength}
  A real-valued physical length, allowing signed optical distances
\end{definition}
\begin{proof}
  This is the declared model definition of Optical Length.
\end{proof}
\begin{definition}[unit Choices For Length]
  \label{def:physics:phyx-mini-0112:phyxminiproblems-problemphyxmini0112-unitchoicesforlength}
  \lean{PhyXMiniProblems.ProblemPhyXMini0112.unitChoicesForLength}
  Unit choices obtained from SI by replacing only the length unit
\end{definition}
\begin{proof}
  This is the declared model definition of unit Choices For Length.
\end{proof}
\begin{definition}[length Readout]
  \label{def:physics:phyx-mini-0112:phyxminiproblems-problemphyxmini0112-lengthreadout}
  \lean{PhyXMiniProblems.ProblemPhyXMini0112.lengthReadout}
  \uses{def:physics:phyx-mini-0112:phyxminiproblems-problemphyxmini0112-opticallength, def:physics:phyx-mini-0112:phyxminiproblems-problemphyxmini0112-unitchoicesforlength}
  Scalar readout of a dimensionful optical length in the chosen length unit
\end{definition}
\begin{proof}
  This is the declared model definition of length Readout.
\end{proof}
\begin{definition}[Lens Label]
  \label{def:physics:phyx-mini-0112:phyxminiproblems-problemphyxmini0112-lenslabel}
  \lean{PhyXMiniProblems.ProblemPhyXMini0112.LensLabel}
  The two lenses in propagation order
\end{definition}
\begin{proof}
  This is the declared model definition of Lens Label.
\end{proof}
\begin{definition}[Thin Lens Kind]
  \label{def:physics:phyx-mini-0112:phyxminiproblems-problemphyxmini0112-thinlenskind}
  \lean{PhyXMiniProblems.ProblemPhyXMini0112.ThinLensKind}
  Qualitative sign role of a thin lens
\end{definition}
\begin{proof}
  This is the declared model definition of Thin Lens Kind.
\end{proof}
\begin{definition}[Ray Bundle Profile]
  \label{def:physics:phyx-mini-0112:phyxminiproblems-problemphyxmini0112-raybundleprofile}
  \lean{PhyXMiniProblems.ProblemPhyXMini0112.RayBundleProfile}
  Ray-bundle profile before entering the first lens
\end{definition}
\begin{proof}
  This is the declared model definition of Ray Bundle Profile.
\end{proof}
\begin{definition}[Optical Approximation]
  \label{def:physics:phyx-mini-0112:phyxminiproblems-problemphyxmini0112-opticalapproximation}
  \lean{PhyXMiniProblems.ProblemPhyXMini0112.OpticalApproximation}
  Optical approximation under which the ray-transfer laws are asserted
\end{definition}
\begin{proof}
  This is the declared model definition of Optical Approximation.
\end{proof}
\begin{definition}[Axis Point]
  \label{def:physics:phyx-mini-0112:phyxminiproblems-problemphyxmini0112-axispoint}
  \lean{PhyXMiniProblems.ProblemPhyXMini0112.AxisPoint}
  Named points on the optical axis in the supplied figure
\end{definition}
\begin{proof}
  This is the declared model definition of Axis Point.
\end{proof}
\begin{definition}[lens Center Point]
  \label{def:physics:phyx-mini-0112:phyxminiproblems-problemphyxmini0112-lenscenterpoint}
  \lean{PhyXMiniProblems.ProblemPhyXMini0112.lensCenterPoint}
  \uses{def:physics:phyx-mini-0112:phyxminiproblems-problemphyxmini0112-lenslabel, def:physics:phyx-mini-0112:phyxminiproblems-problemphyxmini0112-axispoint}
  The center label belonging to each lens in the figure
\end{definition}
\begin{proof}
  This is the declared model definition of lens Center Point.
\end{proof}
\begin{definition}[Thin Lens]
  \label{def:physics:phyx-mini-0112:phyxminiproblems-problemphyxmini0112-thinlens}
  \lean{PhyXMiniProblems.ProblemPhyXMini0112.ThinLens}
  \uses{def:physics:phyx-mini-0112:phyxminiproblems-problemphyxmini0112-opticallength, def:physics:phyx-mini-0112:phyxminiproblems-problemphyxmini0112-thinlenskind}
  A thin lens carrying its qualitative kind and signed focal length
\end{definition}
\begin{proof}
  This is the declared model definition of Thin Lens.
\end{proof}
\begin{definition}[Zoom Lens Diagram]
  \label{def:physics:phyx-mini-0112:phyxminiproblems-problemphyxmini0112-zoomlensdiagram}
  \lean{PhyXMiniProblems.ProblemPhyXMini0112.ZoomLensDiagram}
  \uses{def:physics:phyx-mini-0112:phyxminiproblems-problemphyxmini0112-opticallength, def:physics:phyx-mini-0112:phyxminiproblems-problemphyxmini0112-lenslabel, def:physics:phyx-mini-0112:phyxminiproblems-problemphyxmini0112-raybundleprofile, def:physics:phyx-mini-0112:phyxminiproblems-problemphyxmini0112-opticalapproximation, def:physics:phyx-mini-0112:phyxminiproblems-problemphyxmini0112-axispoint, def:physics:phyx-mini-0112:phyxminiproblems-problemphyxmini0112-thinlens}
  Physical quantities and labels in the zoom-lens diagram. incomingRayRadius is the displayed r₀, radiusAtSecondLens is r'₀, secondLensImageDistance is s'₂, and effectiveFocalLength is f
\end{definition}
\begin{proof}
  This is the declared model definition of Zoom Lens Diagram.
\end{proof}
\begin{definition}[focal Length Readout]
  \label{def:physics:phyx-mini-0112:phyxminiproblems-problemphyxmini0112-focallengthreadout}
  \lean{PhyXMiniProblems.ProblemPhyXMini0112.focalLengthReadout}
  \uses{def:physics:phyx-mini-0112:phyxminiproblems-problemphyxmini0112-lengthreadout, def:physics:phyx-mini-0112:phyxminiproblems-problemphyxmini0112-lenslabel, def:physics:phyx-mini-0112:phyxminiproblems-problemphyxmini0112-zoomlensdiagram}
  Signed focal-length readout for one lens
\end{definition}
\begin{proof}
  This is the declared model definition of focal Length Readout.
\end{proof}
\begin{definition}[axis Position Readout]
  \label{def:physics:phyx-mini-0112:phyxminiproblems-problemphyxmini0112-axispositionreadout}
  \lean{PhyXMiniProblems.ProblemPhyXMini0112.axisPositionReadout}
  \uses{def:physics:phyx-mini-0112:phyxminiproblems-problemphyxmini0112-lengthreadout, def:physics:phyx-mini-0112:phyxminiproblems-problemphyxmini0112-axispoint, def:physics:phyx-mini-0112:phyxminiproblems-problemphyxmini0112-zoomlensdiagram}
  Signed coordinate readout of a labelled point on the optical axis
\end{definition}
\begin{proof}
  This is the declared model definition of axis Position Readout.
\end{proof}
\begin{definition}[axial Separation Readout]
  \label{def:physics:phyx-mini-0112:phyxminiproblems-problemphyxmini0112-axialseparationreadout}
  \lean{PhyXMiniProblems.ProblemPhyXMini0112.axialSeparationReadout}
  \uses{def:physics:phyx-mini-0112:phyxminiproblems-problemphyxmini0112-axispoint, def:physics:phyx-mini-0112:phyxminiproblems-problemphyxmini0112-zoomlensdiagram, def:physics:phyx-mini-0112:phyxminiproblems-problemphyxmini0112-axispositionreadout}
  Signed axial separation from start to finish
\end{definition}
\begin{proof}
  This is the declared model definition of axial Separation Readout.
\end{proof}
\begin{definition}[lens Separation Readout]
  \label{def:physics:phyx-mini-0112:phyxminiproblems-problemphyxmini0112-lensseparationreadout}
  \lean{PhyXMiniProblems.ProblemPhyXMini0112.lensSeparationReadout}
  \uses{def:physics:phyx-mini-0112:phyxminiproblems-problemphyxmini0112-zoomlensdiagram, def:physics:phyx-mini-0112:phyxminiproblems-problemphyxmini0112-axialseparationreadout}
  The adjustable lens separation d
\end{definition}
\begin{proof}
  This is the declared model definition of lens Separation Readout.
\end{proof}
\begin{definition}[incoming Ray Radius Readout]
  \label{def:physics:phyx-mini-0112:phyxminiproblems-problemphyxmini0112-incomingrayradiusreadout}
  \lean{PhyXMiniProblems.ProblemPhyXMini0112.incomingRayRadiusReadout}
  \uses{def:physics:phyx-mini-0112:phyxminiproblems-problemphyxmini0112-lengthreadout, def:physics:phyx-mini-0112:phyxminiproblems-problemphyxmini0112-zoomlensdiagram}
  Readout of the displayed ray radius r₀
\end{definition}
\begin{proof}
  This is the declared model definition of incoming Ray Radius Readout.
\end{proof}
\begin{definition}[radius At Second Lens Readout]
  \label{def:physics:phyx-mini-0112:phyxminiproblems-problemphyxmini0112-radiusatsecondlensreadout}
  \lean{PhyXMiniProblems.ProblemPhyXMini0112.radiusAtSecondLensReadout}
  \uses{def:physics:phyx-mini-0112:phyxminiproblems-problemphyxmini0112-lengthreadout, def:physics:phyx-mini-0112:phyxminiproblems-problemphyxmini0112-zoomlensdiagram}
  Readout of the displayed ray radius r'₀ at the second lens
\end{definition}
\begin{proof}
  This is the declared model definition of radius At Second Lens Readout.
\end{proof}
\begin{definition}[second Lens Image Distance Readout]
  \label{def:physics:phyx-mini-0112:phyxminiproblems-problemphyxmini0112-secondlensimagedistancereadout}
  \lean{PhyXMiniProblems.ProblemPhyXMini0112.secondLensImageDistanceReadout}
  \uses{def:physics:phyx-mini-0112:phyxminiproblems-problemphyxmini0112-lengthreadout, def:physics:phyx-mini-0112:phyxminiproblems-problemphyxmini0112-zoomlensdiagram}
  Readout of the second-lens image distance s'₂
\end{definition}
\begin{proof}
  This is the declared model definition of second Lens Image Distance Readout.
\end{proof}
\begin{definition}[effective Focal Length Readout]
  \label{def:physics:phyx-mini-0112:phyxminiproblems-problemphyxmini0112-effectivefocallengthreadout}
  \lean{PhyXMiniProblems.ProblemPhyXMini0112.effectiveFocalLengthReadout}
  \uses{def:physics:phyx-mini-0112:phyxminiproblems-problemphyxmini0112-lengthreadout, def:physics:phyx-mini-0112:phyxminiproblems-problemphyxmini0112-zoomlensdiagram}
  Readout of the effective focal length f of the two-lens combination
\end{definition}
\begin{proof}
  This is the declared model definition of effective Focal Length Readout.
\end{proof}
\begin{definition}[Matches Zoom Lens Figure]
  \label{def:physics:phyx-mini-0112:phyxminiproblems-problemphyxmini0112-matcheszoomlensfigure}
  \lean{PhyXMiniProblems.ProblemPhyXMini0112.MatchesZoomLensFigure}
  \uses{def:physics:phyx-mini-0112:phyxminiproblems-problemphyxmini0112-zoomlensdiagram, def:physics:phyx-mini-0112:phyxminiproblems-problemphyxmini0112-axialseparationreadout, def:physics:phyx-mini-0112:phyxminiproblems-problemphyxmini0112-secondlensimagedistancereadout, def:physics:phyx-mini-0112:phyxminiproblems-problemphyxmini0112-effectivefocallengthreadout}
  Qualitative labels and distance annotations read from the supplied figure. The predicate identifies s'₂ and f with their drawn axial spans but gives no numerical value for the requested separation
\end{definition}
\begin{proof}
  This is the declared model definition of Matches Zoom Lens Figure.
\end{proof}
\begin{definition}[Matches Zoom Lens Problem Data]
  \label{def:physics:phyx-mini-0112:phyxminiproblems-problemphyxmini0112-matcheszoomlensproblemdata}
  \lean{PhyXMiniProblems.ProblemPhyXMini0112.MatchesZoomLensProblemData}
  \uses{def:physics:phyx-mini-0112:phyxminiproblems-problemphyxmini0112-zoomlensdiagram, def:physics:phyx-mini-0112:phyxminiproblems-problemphyxmini0112-focallengthreadout, def:physics:phyx-mini-0112:phyxminiproblems-problemphyxmini0112-lensseparationreadout, def:physics:phyx-mini-0112:phyxminiproblems-problemphyxmini0112-effectivefocallengthreadout}
  Numerical problem data in centimeters. The desired effective focal length is an input condition of the question; the separation itself is constrained only to the stated adjustable interval
\end{definition}
\begin{proof}
  This is the declared model definition of Matches Zoom Lens Problem Data.
\end{proof}
\begin{definition}[Has Physical Zoom Lens Configuration]
  \label{def:physics:phyx-mini-0112:phyxminiproblems-problemphyxmini0112-hasphysicalzoomlensconfiguration}
  \lean{PhyXMiniProblems.ProblemPhyXMini0112.HasPhysicalZoomLensConfiguration}
  \uses{def:physics:phyx-mini-0112:phyxminiproblems-problemphyxmini0112-zoomlensdiagram, def:physics:phyx-mini-0112:phyxminiproblems-problemphyxmini0112-focallengthreadout, def:physics:phyx-mini-0112:phyxminiproblems-problemphyxmini0112-axispositionreadout, def:physics:phyx-mini-0112:phyxminiproblems-problemphyxmini0112-lensseparationreadout, def:physics:phyx-mini-0112:phyxminiproblems-problemphyxmini0112-incomingrayradiusreadout, def:physics:phyx-mini-0112:phyxminiproblems-problemphyxmini0112-radiusatsecondlensreadout, def:physics:phyx-mini-0112:phyxminiproblems-problemphyxmini0112-secondlensimagedistancereadout, def:physics:phyx-mini-0112:phyxminiproblems-problemphyxmini0112-effectivefocallengthreadout}
  Sign, ordering, and nondegeneracy conditions for the physical branch shown in the diagram. The final inequality is the source condition |f₂| > f₁ - d
\end{definition}
\begin{proof}
  This is the declared model definition of Has Physical Zoom Lens Configuration.
\end{proof}
\begin{definition}[Satisfies Paraxial Zoom Lens Laws]
  \label{def:physics:phyx-mini-0112:phyxminiproblems-problemphyxmini0112-satisfiesparaxialzoomlenslaws}
  \lean{PhyXMiniProblems.ProblemPhyXMini0112.SatisfiesParaxialZoomLensLaws}
  \uses{def:physics:phyx-mini-0112:phyxminiproblems-problemphyxmini0112-zoomlensdiagram, def:physics:phyx-mini-0112:phyxminiproblems-problemphyxmini0112-focallengthreadout, def:physics:phyx-mini-0112:phyxminiproblems-problemphyxmini0112-lensseparationreadout, def:physics:phyx-mini-0112:phyxminiproblems-problemphyxmini0112-incomingrayradiusreadout, def:physics:phyx-mini-0112:phyxminiproblems-problemphyxmini0112-radiusatsecondlensreadout, def:physics:phyx-mini-0112:phyxminiproblems-problemphyxmini0112-secondlensimagedistancereadout, def:physics:phyx-mini-0112:phyxminiproblems-problemphyxmini0112-effectivefocallengthreadout}
  Paraxial ray-transfer and thin-lens laws for the separated lens pair. For every length unit, the first equation transports the ray height from r₀ to r'₀; the second adds the angular deflections of the two thin lenses; the third is the signed Gaussian equation for the diverging lens, whose virtual object distance is d - f₁. All equations are division-free and homogeneous in physical dimension
\end{definition}
\begin{proof}
  This is the declared model definition of Satisfies Paraxial Zoom Lens Laws.
\end{proof}
\begin{lemma}[separated lens effective focal relation]
  \label{lem:physics:phyx-mini-0112:phyxminiproblems-problemphyxmini0112-separated-lens-effective-focal-relation}
  \lean{PhyXMiniProblems.ProblemPhyXMini0112.separated_lens_effective_focal_relation}
  \uses{def:physics:phyx-mini-0112:phyxminiproblems-problemphyxmini0112-zoomlensdiagram, def:physics:phyx-mini-0112:phyxminiproblems-problemphyxmini0112-focallengthreadout, def:physics:phyx-mini-0112:phyxminiproblems-problemphyxmini0112-lensseparationreadout, def:physics:phyx-mini-0112:phyxminiproblems-problemphyxmini0112-effectivefocallengthreadout, def:physics:phyx-mini-0112:phyxminiproblems-problemphyxmini0112-hasphysicalzoomlensconfiguration, def:physics:phyx-mini-0112:phyxminiproblems-problemphyxmini0112-satisfiesparaxialzoomlenslaws}
  The ray-height and refraction laws imply the standard effective-focal-length relation for two separated thin lenses, f (f₁ + f₂ - d) = f₁ f₂
\end{lemma}
\begin{proof}
  The conclusion follows from the governing relations and figure readouts named in the statement of separated lens effective focal relation.
\end{proof}
\begin{definition}[Answer Choice]
  \label{def:physics:phyx-mini-0112:phyxminiproblems-problemphyxmini0112-answerchoice}
  \lean{PhyXMiniProblems.ProblemPhyXMini0112.AnswerChoice}
  Labels of the four displayed answer choices
\end{definition}
\begin{proof}
  This is the declared model definition of Answer Choice.
\end{proof}
\begin{definition}[separation In Centimeters]
  \label{def:physics:phyx-mini-0112:phyxminiproblems-problemphyxmini0112-answerchoice-separationincentimeters}
  \lean{PhyXMiniProblems.ProblemPhyXMini0112.AnswerChoice.separationInCentimeters}
  \uses{def:physics:phyx-mini-0112:phyxminiproblems-problemphyxmini0112-answerchoice}
  Separation in centimeters printed beside each answer choice
\end{definition}
\begin{proof}
  This is the declared model definition of separation In Centimeters.
\end{proof}
\begin{definition}[recorded Answer Choice]
  \label{def:physics:phyx-mini-0112:phyxminiproblems-problemphyxmini0112-recordedanswerchoice}
  \lean{PhyXMiniProblems.ProblemPhyXMini0112.recordedAnswerChoice}
  \uses{def:physics:phyx-mini-0112:phyxminiproblems-problemphyxmini0112-answerchoice}
  Dataset metadata: the recorded answer label, not a theorem premise
\end{definition}
\begin{proof}
  This is the declared model definition of recorded Answer Choice.
\end{proof}
\begin{definition}[Matches Answer]
  \label{def:physics:phyx-mini-0112:phyxminiproblems-problemphyxmini0112-matchesanswer}
  \lean{PhyXMiniProblems.ProblemPhyXMini0112.MatchesAnswer}
  \uses{def:physics:phyx-mini-0112:phyxminiproblems-problemphyxmini0112-zoomlensdiagram, def:physics:phyx-mini-0112:phyxminiproblems-problemphyxmini0112-lensseparationreadout, def:physics:phyx-mini-0112:phyxminiproblems-problemphyxmini0112-answerchoice}
  Exact agreement between the modeled lens separation and an answer choice
\end{definition}
\begin{proof}
  This is the declared model definition of Matches Answer.
\end{proof}
% --- Archon declaration topology end ---
% --- Archon physics formalization source end ---
